\documentclass[aspectratio=169,12pt]{beamer}
% Estilo canónico para las presentaciones de Métodos Numéricos 2026-2.
\usepackage[T1]{fontenc}
\usepackage[utf8]{inputenc}
\usepackage[spanish,es-nodecimaldot]{babel}
\usepackage{lmodern}
\usepackage{amsmath,amssymb,mathtools}
\usepackage{booktabs}
\usepackage{graphicx}
\usepackage{xcolor}
\usepackage{hyperref}

\definecolor{upazul}{HTML}{003B70}
\definecolor{upazulclaro}{HTML}{EAF2F8}
\definecolor{updorado}{HTML}{C5962E}
\definecolor{upgris}{HTML}{4A5560}
\definecolor{upverde}{HTML}{167C5A}
\definecolor{uprojo}{HTML}{B33A3A}

\setbeamercolor{normal text}{fg=black,bg=white}
\setbeamercolor{structure}{fg=upazul}
\setbeamercolor{frametitle}{fg=white,bg=upazul}
\setbeamercolor{title}{fg=upazul}
\setbeamercolor{subtitle}{fg=upgris}
\setbeamercolor{block title}{fg=white,bg=upazul}
\setbeamercolor{block body}{fg=black,bg=upazulclaro}
\setbeamercolor{alerted text}{fg=uprojo}
\setbeamercolor{example text}{fg=upverde}

\setbeamerfont{title}{size=\Huge,series=\bfseries}
\setbeamerfont{subtitle}{size=\Large}
\setbeamerfont{frametitle}{size=\LARGE,series=\bfseries}
\setbeamerfont{normal text}{size=\large}
\setbeamerfont{block title}{series=\bfseries}

\setbeamertemplate{navigation symbols}{}
\setbeamertemplate{itemize item}{\color{updorado}$\blacktriangleright$}
\setbeamertemplate{itemize subitem}{\color{upazul}$\bullet$}
\setbeamertemplate{blocks}[rounded][shadow=false]

\setbeamertemplate{frametitle}{%
  \nointerlineskip
  \begin{beamercolorbox}[wd=\paperwidth,ht=0.78cm,dp=0.24cm,leftskip=0.55cm]{frametitle}
    \usebeamerfont{frametitle}\insertframetitle
  \end{beamercolorbox}}

\setbeamertemplate{footline}{%
  \leavevmode
  \hbox{%
    \begin{beamercolorbox}[wd=.82\paperwidth,ht=2.7ex,dp=1.2ex,leftskip=.5cm]{author in head/foot}
      \color{upgris}\scriptsize Métodos Numéricos 2026-2 \hspace{0.8em} | \hspace{0.8em} Universidad del Pacífico
    \end{beamercolorbox}%
    \begin{beamercolorbox}[wd=.18\paperwidth,ht=2.7ex,dp=1.2ex,rightskip=.5cm plus1fil]{date in head/foot}
      \color{upazul}\scriptsize\hfill\insertframenumber/\inserttotalframenumber
    \end{beamercolorbox}}
  \vskip0pt}

\newcommand{\concepto}[1]{\textcolor{upazul}{\textbf{#1}}}
\newcommand{\resultado}[1]{\textcolor{upverde}{\textbf{#1}}}
\newcommand{\advertencia}[1]{\textcolor{uprojo}{\textbf{#1}}}
\newcommand{\codigo}[1]{\colorbox{upazulclaro}{\strut\texttt{#1}}}
\newcommand{\fuente}[1]{\vfill\hfill{\tiny\color{upgris}#1}}



\usepackage{array}
\usepackage{tikz}
\usetikzlibrary{arrows.meta,decorations.pathreplacing}

\title{Ecuaciones no lineales}
\subtitle{Bisección, punto fijo, Newton y secante}
\author{Oswaldo José Velásquez Castañón}
\institute{Universidad del Pacífico}
\date{Métodos Numéricos -- 2026-2}

\newcommand{\R}{\mathbb{R}}
\newcommand{\absval}[1]{\left|#1\right|}

\begin{document}

\begin{frame}[plain]
  \begin{columns}[T,onlytextwidth]
    \begin{column}{0.70\textwidth}
      \vspace{0.7cm}
      {\usebeamerfont{title}\color{upazul}\inserttitle\par}
      \vspace{0.35cm}
      {\usebeamerfont{subtitle}\color{upgris}\insertsubtitle\par}
      \vspace{1.1cm}
      {\large\insertauthor\par}
      {\normalsize\insertdate\par}
    \end{column}
    \begin{column}{0.25\textwidth}
      \vspace{0.55cm}
      \includegraphics[width=\linewidth]{UPlogo.pdf}
    \end{column}
  \end{columns}
\end{frame}

\begin{frame}{El problema escalar}
Dada una función continua $f:\R\to\R$, buscamos una raíz $x^*$ que satisfaga
\[
f(x^*)=0.
\]
\begin{block}{Dos decisiones antes de iterar}
Primero localizamos una región plausible. Después elegimos entre un método que
\concepto{garantiza avance} y uno que privilegia la \concepto{velocidad local}.
\end{block}
La función
\[
f(x)=x^3-x-2
\]
servirá como ejemplo común. Su raíz positiva es $x^*\approx1.521379707$.
\end{frame}

\begin{frame}{Una gráfica orienta, pero no certifica}
Un cambio de signo permite construir una garantía matemática.
\begin{block}{Teorema del valor intermedio}
Si $f$ es continua en $[a,b]$ y $f(a)f(b)<0$, entonces existe al menos una raíz
$x^*\in(a,b)$.
\end{block}
Para $f(x)=x^3-x-2$,
\[
f(1)=-2,
\qquad
f(2)=4.
\]
Por tanto, hay una raíz en $(1,2)$. La condición no afirma que la raíz sea única.
\end{frame}

\begin{frame}{Intervalo de bisección}
\begin{columns}[T,onlytextwidth]
\begin{column}{0.46\textwidth}
Partimos de $[a_k,b_k]$ y calculamos
\[
m_k=\frac{a_k+b_k}{2}.
\]
Conservamos la mitad donde los valores extremos tienen signos opuestos.
\begin{block}{Invariante}
$f(a_k)f(b_k)\leq0$ en cada iteración.
\end{block}
\end{column}
\begin{column}{0.52\textwidth}
\centering
\resizebox{\linewidth}{!}{\begin{tikzpicture}[x=1.25cm,y=1.05cm,>=Latex]
  \draw[->] (-0.2,0) -- (6.4,0) node[right] {$x$};
  \draw[->] (0,-1.5) -- (0,2.1) node[above] {$f(x)$};
  \draw[thick,blue!70!black,smooth] plot coordinates
    {(0.4,-1.1) (1.2,-0.8) (2,-0.25) (2.5,0.15) (3.4,0.9) (4.5,1.45) (5.8,1.6)};
  \coordinate (a) at (0.7,0);
  \coordinate (b) at (5.5,0);
  \coordinate (m) at (3.1,0);
  \draw[dashed] (a) -- (0.7,-0.95);
  \draw[dashed] (b) -- (5.5,1.55);
  \draw[dashed] (m) -- (3.1,0.68);
  \fill (a) circle (1.5pt) node[above left] {$a^{(0)}$};
  \fill (b) circle (1.5pt) node[below right] {$b^{(0)}$};
  \fill (m) circle (1.5pt) node[below] {$x^{(1)}$};
  \fill (2.3,0) circle (1.7pt) node[below left] {$x^*$};
  \draw[decorate,decoration={brace,mirror,amplitude=4pt}] (a) -- (m)
    node[midway,below=7pt] {nuevo intervalo};
\end{tikzpicture}
}
\end{column}
\end{columns}
\end{frame}

\begin{frame}{Algoritmo de bisección}
\begin{block}{Entrada}
$f$ continua, $f(a)f(b)<0$, tolerancia $\varepsilon>0$ y máximo de iteraciones.
\end{block}
\begin{enumerate}
  \item Calcular $m=(a+b)/2$.
  \item Si $f(m)=0$, terminar.
  \item Si $f(a)f(m)<0$, reemplazar $b$ por $m$. En otro caso, reemplazar $a$ por $m$.
  \item Repetir hasta que $(b-a)/2\leq\varepsilon$ o el residuo sea aceptable.
\end{enumerate}
\begin{block}{Salida certificada}
El punto medio final satisface $\absval{m-x^*}\leq(b-a)/2$.
\end{block}
\end{frame}

\begin{frame}{Bisección ofrece una cota antes de calcular}
Después de $k$ reducciones,
\[
b_k-a_k=\frac{b_0-a_0}{2^k},
\qquad
\absval{m_k-x^*}\leq\frac{b_0-a_0}{2^{k+1}}.
\]
Para garantizar $\absval{m_k-x^*}\leq\varepsilon$, basta tomar
\[
k\geq
\left\lceil
\log_2\!\left(\frac{b_0-a_0}{2\varepsilon}\right)
\right\rceil.
\]
\begin{block}{Lectura}
Cada iteración gana un bit de precisión. La convergencia es lineal y predecible.
\end{block}
\end{frame}

\begin{frame}{Bisección aplicada a $x^3-x-2=0$}
\centering
\begin{tabular}{ccccc}
\toprule
$k$ & $a_k$ & $b_k$ & $m_k$ & $f(m_k)$ \\
\midrule
0 & 1.000000 & 2.000000 & 1.500000 & $-0.125000$ \\
1 & 1.500000 & 2.000000 & 1.750000 & $ 1.609375$ \\
2 & 1.500000 & 1.750000 & 1.625000 & $ 0.666016$ \\
3 & 1.500000 & 1.625000 & 1.562500 & $ 0.252197$ \\
4 & 1.500000 & 1.562500 & 1.531250 & $ 0.059113$ \\
5 & 1.500000 & 1.531250 & 1.515625 & $-0.034054$ \\
\bottomrule
\end{tabular}
\vspace{0.35cm}

El método nunca pierde la raíz encerrada, aunque el punto medio cambie de lado.
\end{frame}

\begin{frame}{El criterio de parada depende del objetivo}
En una implementación conviene vigilar dos cantidades:
\[
\text{paso relativo}
=\frac{|x_{k+1}-x_k|}{\max\{1,|x_{k+1}|\}},
\qquad
\text{residuo}=|f(x_{k+1})|.
\]
\begin{itemize}
  \item Un paso pequeño indica estancamiento, pero no certifica una raíz.
  \item Un residuo pequeño indica que la ecuación casi se satisface.
  \item En bisección, el ancho del intervalo aporta una cota directa del error.
\end{itemize}
\advertencia{Siempre se fija también un máximo de iteraciones.}
\end{frame}

\begin{frame}{Punto fijo cambia la pregunta}
Una ecuación $f(x)=0$ puede reformularse como
\[
x=g(x).
\]
A partir de $x_0$, la iteración toma la forma
\[
x_{k+1}=g(x_k).
\]
\begin{block}{La elección de $g$ importa}
Dos reformulaciones algebraicamente equivalentes pueden producir una sucesión
convergente y otra divergente.
\end{block}
Ejemplo: $x=\cos x$ genera $x_{k+1}=\cos x_k$.
\end{frame}

\begin{frame}{El diagrama de telaraña muestra la dinámica}
\begin{columns}[T,onlytextwidth]
\begin{column}{0.43\textwidth}
La trayectoria alterna entre:
\begin{itemize}
  \item la curva $y=g(x)$, donde se evalúa $x_{k+1}$;
  \item la recta $y=x$, donde el nuevo valor pasa al eje horizontal.
\end{itemize}
La intersección es el punto fijo $x^*$.
\end{column}
\begin{column}{0.55\textwidth}
\centering
\resizebox{0.86\linewidth}{!}{\begin{tikzpicture}[x=1.15cm,y=1.15cm,>=Latex]
  \draw[->] (-0.2,0) -- (5.6,0) node[right] {$x$};
  \draw[->] (0,-0.2) -- (0,5.4) node[above] {$y$};
  \draw[gray,thick] (0.3,0.3) -- (5,5) node[above] {$y=x$};
  \draw[blue!70!black,thick,smooth] plot coordinates
    {(0.3,4.5) (0.8,3.8) (1.4,3.15) (2,2.6) (2.6,2.15) (3.3,1.75) (4.1,1.42) (5,1.18)}
    node[right] {$y=F(x)$};
  \draw[orange!80!black,-{Latex[length=2mm]}] (0.8,0) -- (0.8,3.8) -- (3.8,3.8) -- (3.8,1.53) -- (1.53,1.53) -- (1.53,3.02) -- (3.02,3.02);
  \node[below] at (0.8,0) {$x^{(0)}$};
  \node[below] at (3.8,0) {$x^{(1)}$};
  \fill (2.35,2.35) circle (1.6pt) node[below right] {$x^*$};
\end{tikzpicture}
}
\end{column}
\end{columns}
\end{frame}

\begin{frame}{Una contracción garantiza convergencia}
Sea $g:[a,b]\to[a,b]$. Si existe $L<1$ tal que
\[
|g(x)-g(y)|\leq L|x-y|
\qquad\text{para todo }x,y\in[a,b],
\]
entonces $g$ tiene un único punto fijo y $x_{k+1}=g(x_k)$ converge a él para
cualquier $x_0\in[a,b]$.
\begin{block}{Cota a posteriori}
\[
|x_k-x^*|\leq\frac{L}{1-L}|x_k-x_{k-1}|.
\]
\end{block}
La cota traduce un paso observable en una estimación del error.
\end{frame}

\begin{frame}{La derivada permite verificar la contracción}
Si $g$ es derivable y
\[
\max_{x\in[a,b]}|g'(x)|\leq L<1,
\]
entonces $g$ es una contracción en $[a,b]$.
\begin{itemize}
  \item $|g'(x^*)|<1$ anticipa convergencia local.
  \item Cuanto menor sea $|g'(x^*)|$, más rápida será la convergencia lineal.
  \item $|g'(x^*)|>1$ hace que el punto fijo sea localmente repulsivo.
\end{itemize}
\begin{block}{Diseño de la iteración}
La reformulación debe mantener las iteraciones en el intervalo y reducir las distancias.
\end{block}
\end{frame}

\begin{frame}{Punto fijo para $x=\cos x$}
Con $x_0=1$,
\[
x_{k+1}=\cos x_k.
\]
\centering
\begin{tabular}{ccc}
\toprule
$k$ & $x_k$ & $|x_{k+1}-x_k|$ \\
\midrule
0 & 1.000000 & 0.459698 \\
1 & 0.540302 & 0.317251 \\
2 & 0.857553 & 0.203263 \\
3 & 0.654290 & 0.139191 \\
4 & 0.793480 & 0.092111 \\
5 & 0.701369 & 0.062591 \\
\bottomrule
\end{tabular}

La sucesión oscila y converge linealmente a $x^*\approx0.739085$.
\end{frame}

\begin{frame}{Newton usa una aproximación lineal local}
\begin{columns}[T,onlytextwidth]
\begin{column}{0.45\textwidth}
La expansión de Taylor en $x_k$ da
\[
f(x_k+s)\approx f(x_k)+f'(x_k)s.
\]
Al imponer que el modelo lineal sea cero,
\[
s_k=-\frac{f(x_k)}{f'(x_k)},
\qquad
\boxed{x_{k+1}=x_k+s_k}.
\]
\end{column}
\begin{column}{0.53\textwidth}
\centering
\resizebox{\linewidth}{!}{\begin{tikzpicture}[x=1.2cm,y=1cm,>=Latex]
  \draw[->] (-0.2,0) -- (6.4,0) node[right] {$x$};
  \draw[->] (0,-1.2) -- (0,2.5) node[above] {$f(x)$};
  \draw[thick,blue!70!black,smooth] plot coordinates
    {(0.5,2.0) (1.5,1.35) (2.4,0.45) (2.85,0) (3.5,-0.5) (4.5,-0.55) (5.8,0.6)};
  \coordinate (P) at (5.1,-0.28);
  \draw[orange!80!black,thick] (2.9,0) -- (6.0,-0.62);
  \fill (P) circle (1.7pt) node[below right] {$(x^{(j)},f(x^{(j)}))$};
  \draw[dashed] (5.1,0) -- (P);
  \fill (2.9,0) circle (1.7pt) node[above left] {$x^{(j+1)}$};
  \node[orange!80!black] at (4.5,0.45) {tangente};
\end{tikzpicture}
}
\end{column}
\end{columns}
\end{frame}

\begin{frame}{Algoritmo de Newton}
\begin{block}{Entrada}
$f$, su derivada $f'$, un valor inicial $x_0$, tolerancias y un máximo de iteraciones.
\end{block}
Para $k=0,1,\ldots$:
\begin{enumerate}
  \item evaluar $f(x_k)$ y $f'(x_k)$;
  \item detenerse o cambiar de estrategia si $|f'(x_k)|$ es demasiado pequeño;
  \item calcular $s_k=-f(x_k)/f'(x_k)$ y $x_{k+1}=x_k+s_k$;
  \item comprobar paso y residuo.
\end{enumerate}
\advertencia{Una fórmula correcta no compensa un punto inicial inadecuado.}
\end{frame}

\begin{frame}{Newton aplicado a $x^3-x-2=0$}
Con $x_0=1.5$,
\[
x_{k+1}=x_k-\frac{x_k^3-x_k-2}{3x_k^2-1}.
\]
\centering
\begin{tabular}{cccc}
\toprule
$k$ & $x_k$ & $|f(x_k)|$ & $|x_k-x^*|$ \\
\midrule
0 & 1.500000000 & $1.25\times10^{-1}$ & $2.14\times10^{-2}$ \\
1 & 1.521739130 & $2.14\times10^{-3}$ & $3.59\times10^{-4}$ \\
2 & 1.521379806 & $5.89\times10^{-7}$ & $9.92\times10^{-8}$ \\
3 & 1.521379707 & $4.53\times10^{-14}$ & $7.55\times10^{-15}$ \\
\bottomrule
\end{tabular}

El número de cifras correctas casi se duplica cerca de la raíz.
\end{frame}

\begin{frame}{La convergencia cuadrática es local}
Si $f(x^*)=0$, $f'(x^*)\neq0$ y $f''$ permanece acotada cerca de $x^*$, entonces
Newton satisface, para $x_k$ suficientemente cercano,
\[
|e_{k+1}|\leq C|e_k|^2,
\qquad e_k=x_k-x^*.
\]
\begin{block}{Consecuencia}
Un error pequeño se eleva al cuadrado en cada paso. Lejos de la raíz, Newton
puede alejarse, ciclar o encontrar otra raíz.
\end{block}
La teoría combina una condición local con una elección razonable de $x_0$.
\end{frame}

\begin{frame}{Las raíces múltiples frenan a Newton}
Para $f(x)=(x-1)^2$, Newton produce
\[
x_{k+1}=x_k-\frac{(x_k-1)^2}{2(x_k-1)}
=\frac{x_k+1}{2}.
\]
Por tanto,
\[
|x_{k+1}-1|=\frac12|x_k-1|,
\]
y la convergencia pasa de cuadrática a lineal.
\begin{block}{Corrección si se conoce la multiplicidad $m$}
\[
x_{k+1}=x_k-m\frac{f(x_k)}{f'(x_k)}.
\]
\end{block}
\end{frame}

\begin{frame}{Método de la secante}
\begin{columns}[T,onlytextwidth]
\begin{column}{0.48\textwidth}
Con dos aproximaciones consecutivas,
\[
f'(x_k)\approx
\frac{f(x_k)-f(x_{k-1})}{x_k-x_{k-1}}.
\]
Así obtenemos
\[
x_{k+1}=x_k-f(x_k)
\frac{x_k-x_{k-1}}{f(x_k)-f(x_{k-1})}.
\]
\end{column}
\begin{column}{0.50\textwidth}
\centering
\resizebox{\linewidth}{!}{\begin{tikzpicture}[x=1.15cm,y=1cm,>=Latex]
  \draw[->] (-0.2,0) -- (6.5,0) node[right] {$x$};
  \draw[->] (0,-1.4) -- (0,2.5) node[above] {$f(x)$};
  \draw[thick,blue!70!black,smooth] plot coordinates
    {(0.4,2.0) (1.2,1.25) (2.2,0.3) (2.7,-0.15) (3.5,-0.55) (4.5,-0.35) (5.8,0.9)};
  \coordinate (P0) at (1.1,1.35);
  \coordinate (P1) at (4.7,-0.22);
  \draw[orange!80!black,thick] (P0) -- (5.4,-0.52);
  \fill (P0) circle (1.7pt) node[above] {$(x^{(j-1)},f(x^{(j-1)}))$};
  \fill (P1) circle (1.7pt) node[below] {$(x^{(j)},f(x^{(j)}))$};
  \fill (4.20,0) circle (1.7pt) node[above right] {$x^{(j+1)}$};
  \draw[dashed] (1.1,0)--(P0) (4.7,0)--(P1);
\end{tikzpicture}
}
\end{column}
\end{columns}
\end{frame}

\begin{frame}{Algoritmo de la secante}
\begin{block}{Entrada}
Dos valores iniciales $x_0,x_1$, tolerancias y máximo de iteraciones.
\end{block}
Para $k=1,2,\ldots$:
\begin{enumerate}
  \item calcular $d_k=f(x_k)-f(x_{k-1})$;
  \item detenerse si $|d_k|$ es demasiado pequeño;
  \item actualizar $x_{k+1}=x_k-f(x_k)(x_k-x_{k-1})/d_k$;
  \item comprobar paso y residuo.
\end{enumerate}
Cada paso nuevo requiere una sola evaluación adicional de $f$ si se reutilizan
los valores anteriores.
\end{frame}

\begin{frame}{Secante aplicada a $x^3-x-2=0$}
Con $x_0=1$ y $x_1=2$:
\centering
\begin{tabular}{ccc}
\toprule
$k$ & $x_k$ & $|f(x_k)|$ \\
\midrule
0 & 1.000000000 & $2.000000$ \\
1 & 2.000000000 & $4.000000$ \\
2 & 1.333333333 & $0.962963$ \\
3 & 1.462686567 & $0.333339$ \\
4 & 1.531169432 & $0.058626$ \\
5 & 1.520926421 & $0.002693$ \\
6 & 1.521376317 & $2.02\times10^{-5}$ \\
\bottomrule
\end{tabular}

Su orden teórico es $\varphi=(1+\sqrt5)/2\approx1.618$ para una raíz simple.
\end{frame}

\begin{frame}{Costo, garantía y velocidad}
\centering
\small
\begin{tabular}{p{2.2cm}p{2.5cm}p{2.6cm}p{2.6cm}}
\toprule
Método & Datos iniciales & Ventaja principal & Riesgo principal \\
\midrule
Bisección & Intervalo con cambio de signo & Garantía y cota del error & Convergencia lineal \\
Punto fijo & Una función $g$ y $x_0$ & Iteración muy simple & Depende de la reformulación \\
Newton & $x_0$ y $f'$ & Convergencia cuadrática local & Derivada pequeña o mal inicio \\
Secante & $x_0,x_1$ & Evita calcular $f'$ & No conserva un intervalo \\
\bottomrule
\end{tabular}
\normalsize
\vspace{0.35cm}

La mejor elección depende de la información disponible y de la garantía requerida.
\end{frame}

\begin{frame}{Método híbrido}
Una estrategia práctica mantiene un intervalo $[a_k,b_k]$ con cambio de signo.
\begin{enumerate}
  \item Proponer un paso de Newton o secante.
  \item Aceptarlo si permanece en el intervalo y reduce el residuo.
  \item Usar bisección cuando la propuesta no es segura.
  \item Actualizar el intervalo con el nuevo punto.
\end{enumerate}
\begin{block}{Resultado}
La fase inicial conserva la robustez de bisección. Cerca de la raíz, los pasos
abiertos aportan mayor velocidad.
\end{block}
\end{frame}

\begin{frame}{Actividad de clase}
Para
\[
f(x)=e^{-x}-x,
\qquad x\in[0,1],
\]
realice lo siguiente:
\begin{enumerate}
  \item justifique que existe una raíz en $[0,1]$;
  \item ejecute cuatro pasos de bisección;
  \item ejecute cuatro pasos de punto fijo con $g(x)=e^{-x}$ y $x_0=0$;
  \item ejecute tres pasos de Newton desde $x_0=0.5$;
  \item compare los residuos y el número de evaluaciones.
\end{enumerate}
\end{frame}

\begin{frame}{Ideas centrales}
\begin{itemize}
  \item El cambio de signo convierte continuidad en una garantía computable.
  \item Una iteración de punto fijo converge si la reformulación contrae distancias.
  \item Newton obtiene velocidad cuadrática cerca de una raíz simple.
  \item Secante conserva buena velocidad sin evaluar derivadas.
  \item Paso, residuo y cota de error responden preguntas distintas.
\end{itemize}
\end{frame}

\end{document}
